\documentclass{article}
\usepackage{graphicx}
\usepackage{geometry}
\usepackage{hyperref}
\usepackage{enumitem}

\geometry{a4paper, margin=1in}

\title{Introduction to Data Engineering \\ \large{Lecture Notes - Course 1, Week 1}}
\author{Based on materials by DeepLearning.AI and the textbook "Fundamentals of Data Engineering" by Joe Reis & Matt Housley}
\date{\today}

\begin{document}

\maketitle

\section{Introduction and Overview}
Welcome to the study of Data Engineering. This professional certificate program is designed to provide a comprehensive understanding of the data engineering field, preparing you for a successful career.

\textbf{Program Structure:}
\begin{itemize}
    \item \textbf{Course 1:} Introduction to Data Engineering
    \item \textbf{Course 2:} Source Systems, Data Ingestion, and Pipelines
    \item \textbf{Course 3:} Data Storage and Queries
    \item \textbf{Course 4:} Data Modeling, Transformation, and Serving
\end{itemize}

This week, we will take a high-level look at the field of data engineering. We will cover its lifecycle, history, key stakeholders, the importance of business value, and the critical process of translating stakeholder needs into concrete system requirements.

\section{Data Engineering Defined}

\subsection{What is Data Engineering?}
Initially, data generated by software applications was often seen as a mere byproduct or "exhaust," useful mainly for troubleshooting or monitoring application health. However, as organizations began to recognize the intrinsic value of this data, a new discipline emerged. Data engineering evolved from software engineering to specifically handle the growing volume and variety of data.

A formal definition, provided by Reis and Housley in "Fundamentals of Data Engineering," is as follows:
\begin{quote}
    “Data engineering is the development, implementation, and maintenance of systems and processes that take in raw data and produce high-quality, consistent information that supports downstream use cases, such as analysis and machine learning. Data engineering is the intersection of security, data management, DataOps, data architecture, orchestration, and software engineering.”
\end{quote}
In simpler terms, your job as a data engineer is to get raw data from somewhere, turn it into something useful, and then make it available for downstream use cases.

\subsection{The Data Engineering Lifecycle}
The work of a data engineer is structured around the \textbf{Data Engineering Lifecycle}. This framework helps visualize how data moves from its raw state to a valuable asset.

% TODO: Insert diagram of the Data Engineering Lifecycle from the slides.
% \begin{figure}[h!]
% \centering
% \includegraphics[width=0.8\textwidth]{lifecycle_diagram.png}
% \caption{The Data Engineering Lifecycle.}
% \end{figure}

The five major stages, which form the core focus of a data engineer, are:
\begin{itemize}
    \item \textbf{Generation:} Data originates from various \textit{source systems} (e.g., application databases, IoT sensors, logs from software applications). Data engineers are consumers of this data, not typically the owners of these systems.
    \item \textbf{Storage:} Data is persisted for later use. Storage is a foundational layer that underpins all other stages, which is why it is shown spanning the entire middle block of the lifecycle diagram. It can range from traditional databases to large-scale, cloud-based object storage.
    \item \textbf{Ingestion:} Data is moved from source systems into a central storage layer, such as a data lake or data warehouse.
    \item \textbf{Transformation:} Raw data is converted into a useful, high-quality, and structured format. This includes cleaning, normalizing, aggregating, and applying business logic.
    \item \textbf{Serving:} The processed data is made available to consumers. This can be for \textbf{analytics}, \textbf{machine learning}, or \textbf{Reverse ETL}, where data is sent back to source systems to enrich them.
\end{itemize}
The combination of architecture, systems, and processes that moves data through these stages is collectively known as a \textbf{data pipeline}.

\subsection{Lifecycle Undercurrents}
Several critical concepts, or \textbf{undercurrents}, are not discrete stages but rather foundational principles that span the entire lifecycle. They are essential for building robust, secure, and manageable data systems.
\begin{itemize}
    \item \textbf{Security:} Protecting data both at rest (in storage) and in transit (while moving between systems).
    \item \textbf{Data Management:} Encompasses data governance (rules and policies), data quality, and metadata management (data about data).
    \item \textbf{DataOps:} The practice of applying DevOps principles (like automation, monitoring, and CI/CD) to the data lifecycle to improve quality and reduce cycle time.
    \item \textbf{Data Architecture:} The blueprint that defines how data systems are structured and how they interact.
    \item \textbf{Orchestration:} The automation and coordination of complex data workflows (pipelines).
    \item \textbf{Software Engineering:} The application of software development best practices to build reliable and maintainable data systems.
\end{itemize}

\section{A Brief History of Data Engineering}
Understanding the evolution of data engineering helps explain its current state and future direction.
\begin{description}
    \item[1970s] The invention of \textbf{relational databases} and \textbf{SQL} at IBM established the foundation for structured data management.
    \item[1980s] \textbf{Bill Inmon} pioneered the \textbf{data warehouse} concept, separating analytical data from transactional data to support business decision-making.
    \item[1990s] \textbf{Ralph Kimball} advanced data warehousing with \textbf{data modeling} techniques like the star schema for business intelligence (BI). The mainstream adoption of the \textbf{internet} led to an explosion in data generated by web-first companies like Amazon.
    \item[2000s] The \textbf{"Big Data"} era began. The sheer \textbf{volume, velocity, and variety} of data generated by web giants like Google, Yahoo, and Amazon overwhelmed traditional database systems.
    \begin{itemize}
        \item Google's papers on its internal technologies, including \textbf{MapReduce} (2004), provided a new paradigm for processing massive datasets in a distributed manner.
        \item This inspired Yahoo to create \textbf{Hadoop} (2006), an open-source framework that democratized big data processing and gave rise to the "Big Data Engineer."
        \item \textbf{Amazon Web Services (AWS)} launched its public cloud, offering pay-as-you-go access to scalable compute (EC2) and storage (S3), fundamentally changing how systems are built and deployed.
    \end{itemize}
    \item[2010s-Present] The complexity of early big data tools led to a shift towards \textbf{abstract, simplified, and managed services}, primarily on the cloud. The "Modern Data Stack" emerged, emphasizing modular, best-of-breed tools. The role of the data engineer has moved "up the value chain," focusing less on low-level infrastructure management and more on system design, interoperation, and delivering business value. The "Big Data Engineer" is now simply the "Data Engineer."
\end{description}

\section{The Data Engineer and Stakeholders}
Data engineers do not work in isolation. They are central connectors between those who create data and those who consume it.
\begin{itemize}
    \item \textbf{Upstream Stakeholders:} Primarily \textbf{Software Engineers} who build and maintain the source systems that generate data. A data engineer must collaborate with them to understand the data's format, volume, frequency, and any potential disruptions like schema changes.
    \item \textbf{Downstream Stakeholders:} The end-users who derive value from the data. This diverse group includes:
    \begin{itemize}
        \item \textbf{Data Analysts} who create dashboards and reports to analyze business trends.
        \item \textbf{Data Scientists} who build predictive models.
        \item \textbf{Machine Learning Engineers} who productionize those models.
        \item \textbf{Business Leaders} (Executives, Marketing, Sales) who use insights for strategic decision-making.
    \end{itemize}
\end{itemize}
Effective communication is key. For example, when serving a business analyst, a data engineer needs to discuss requirements like query frequency, acceptable data latency (is an hour old okay, or does it need to be real-time?), and precise data definitions (e.g., how is "daily sales" defined in a global company with multiple time zones?).

\section{From Business Goals to Technical Requirements}
\subsection{Delivering Business Value}
The primary objective of a data engineer is to deliver \textbf{business value}. As data warehouse pioneer Bill Inmon advises, "Go to where the money is." Technology is a means to an end, not the end itself. Business value is ultimately determined by your stakeholders and can take many forms:
\begin{itemize}
    \item Increased Revenue
    \item Cost Savings
    \item Improved Operational Efficiency
    \item Enabling the launch of a new data-driven product
\end{itemize}
The value you provide is often based on the \textit{perception} of your stakeholders—whether they believe your work is helping them achieve their goals.

\subsection{The Requirements Gathering Process}
To build a valuable system, a data engineer must translate high-level goals into specific, actionable requirements. This follows a clear hierarchy:
\begin{enumerate}
    \item \textbf{Business Requirements:} The high-level goals of the business (e.g., "increase user retention").
    \item \textbf{Stakeholder Requirements:} The needs of individuals that contribute to those business goals (e.g., "a data scientist needs customer activity data to build a churn prediction model").
    \item \textbf{System Requirements:} The technical specifications for the data system, which are divided into:
    \begin{itemize}
        \item \textbf{Functional Requirements (The "WHAT"):} What the system must do (e.g., "The system must provide daily data updates").
        \item \textbf{Non-Functional Requirements (The "HOW"):} The technical qualities of the system (e.g., "The data must be available by 8 AM UTC with less than 1-hour latency").
    \end{itemize}
\end{enumerate}

\textbf{Key Tactic:} When gathering requirements, the most insightful question you can ask a stakeholder is not "What do you need?" but \textbf{"What action do you plan to take with the data?"} This reveals the underlying goal and ensures you build a system that is truly useful, not just what was initially requested.

\section{Thinking Like a Data Engineer}
A successful data engineer approaches projects with an iterative, structured mindset. While presented as linear steps, this is a continuous, cyclical process.
\begin{enumerate}
    \item \textbf{Identify business goals \& stakeholder needs:} This is the discovery phase. Understand the high-level business objectives and talk to stakeholders to learn about their needs, existing systems, and pain points. Ask what actions they plan to take with the data.
    \item \textbf{Define system requirements:} Translate stakeholder needs into specific, documented functional and non-functional requirements. Confirm your understanding with stakeholders before proceeding.
    \item \textbf{Choose tools \& technologies:} Identify potential tools and perform a cost-benefit analysis. Build a prototype to test your assumptions and validate that your proposed architecture will meet stakeholder needs. This is a critical step to de-risk the project.
    \item \textbf{Build, evaluate, iterate \& evolve:} Build and deploy the production system. Use monitoring and stakeholder feedback to continuously evaluate and improve the system. Data systems must evolve with the business.
\end{enumerate}

\section{Data Engineering in Practice: The Cloud}
This course takes a \textbf{cloud-first} approach, as modern data engineering is predominantly practiced on public cloud platforms like AWS, GCP, and Azure. We will learn to use specific tools on a \textbf{just-in-time} basis as we encounter them in labs.
\begin{itemize}
    \item \textbf{On-Premises:} The company owns and manages its own data centers. This model provides high control but is capital-intensive and less flexible.
    \item \textbf{Cloud:} Resources are rented from a provider on a pay-as-you-go basis. This offers elasticity, scalability, and democratized access to powerful technologies.
    \item \textbf{Hybrid:} A mix of both on-premises and cloud infrastructure.
\end{itemize}

\subsection{Core Cloud Concepts for Data Engineering}
\begin{itemize}
    \item \textbf{IT Resources:} These are the fundamental building blocks.
    \begin{itemize}
        \item \textbf{Compute:} Services that run code, such as virtual machines (\textbf{Amazon EC2}), serverless functions (\textbf{AWS Lambda}), and container services.
        \item \textbf{Storage:} Services that store data. This includes \textbf{Object Storage} (e.g., \textbf{Amazon S3}) for unstructured data, \textbf{Block Storage} (e.g., \textbf{Amazon EBS}) for high-performance database needs, and \textbf{File Storage} (e.g., \textbf{Amazon EFS}).
        \item \textbf{Databases:} Managed services for different database types, like relational databases (\textbf{Amazon RDS}) and data warehouses (\textbf{Amazon Redshift}).
        \item \textbf{Networking:} Services for creating isolated networks in the cloud, like \textbf{Amazon Virtual Private Cloud (VPC)}.
    \end{itemize}
    \item \textbf{Scalability \& Elasticity:} The ability of the cloud to automatically scale resources up or down to match demand, ensuring performance and cost-efficiency.
    \item \textbf{Regions \& Availability Zones (AZs):} The global infrastructure of cloud providers. A \textit{Region} is a geographic area (e.g., US East). Each Region contains multiple, isolated \textit{Availability Zones}, which are clusters of data centers. This structure is designed for high availability and disaster recovery.
\end{itemize}

\subsection{Security: The Shared Responsibility Model}
Security in the cloud is a partnership. The "high-rise apartment" analogy is useful here: the building owner is responsible for the building's security (locks, secure foundation), but the tenant is responsible for locking their own door.
\begin{itemize}
    \item \textbf{AWS is responsible for security OF the cloud:} This includes the physical security of data centers, the hardware, and the underlying hypervisor.
    \item \textbf{You are responsible for security IN the cloud:} This includes securely configuring your services, managing user permissions, setting up firewalls, patching operating systems on your VMs, and encrypting your data.
\end{itemize}

\section{Week 1 Summary}
This week provided a high-level framework for thinking about data engineering. The key takeaways are:
\begin{enumerate}
    \item \textbf{Understand Stakeholder Needs:} Your work must be driven by the needs of your stakeholders and the overarching goals of the business.
    \item \textbf{Translate Needs to Requirements:} Systematically translate abstract needs into concrete functional and non-functional system requirements.
    \item \textbf{Choose Appropriate Tools:} With a clear set of requirements, you can select the right tools and technologies to build a system that delivers tangible value.
\end{enumerate}
By internalizing this framework, a data engineer can move beyond simply implementing technologies and become a key driver of business success.

\end{document}
